%% ma: L11-B07-0004
%% lop: 11
%% chuong: 2
%% bai: 7
%% dang: 11.7.3
%% loai: TN4
%% muc: TH
%% dap_an: "C"
%% nguon: "(NBV)-ĐỀ SỐ 18-MỖI NGÀY 1 ĐỀ THI 2025.pdf (D18, MỖI NGÀY 1 ĐỀ - Đề số 18), Phần I Câu 9"
%% kien_thuc: "Bài 7 (cấp số nhân, bài toán lãi kép)"
%% trang_thai: chinh-thuc
%% ghi_chu: "Đáp án gốc C đúng (giải độc lập: 10.1,07^5 - 10 = 4,0255 ≈ 4,026). Bỏ ảnh trang trí (người giao dịch tại quầy ngân hàng). Đã duyệt 19/10/2026."
\begin{ex}
Một người gửi tiết kiệm $10$ triệu đồng vào một ngân hàng với lãi suất $7\%$ một năm. Biết rằng nếu không rút tiền ra khỏi ngân hàng thì cứ sau mỗi năm, số tiền lãi sẽ được nhập vào vốn ban đầu. Sau $5$ năm mới rút lãi thì người đó thu được số tiền lãi là
\choice{$14{,}026$ triệu đồng}{$50{,}7$ triệu đồng}{\True $4{,}026$ triệu đồng}{$3{,}5$ triệu đồng}
\loigiai{Số tiền cả vốn lẫn lãi sau $5$ năm là $10\cdot(1+0{,}07)^5\approx14{,}026$ (triệu đồng). Số tiền lãi là $14{,}026-10=4{,}026$ (triệu đồng). Chọn C.}
\end{ex}
